\documentclass[11pt]{article}
\usepackage{fontspec}
\usepackage[french,english]{babel}
\usepackage{amsmath}
\usepackage{amssymb}
\usepackage{amsthm}
\usepackage{mathrsfs}
\usepackage{tikz-cd}
\usepackage{enumitem}

\frenchsetup{AutoSpacePunctuation=false}

\newcommand{\printedpage}[1]{\marginpar{\footnotesize p.~#1}}
\newcommand{\msp}[1]{\marginpar{\footnotesize [#1]}}

\newcommand{\M}{\mathcal{M}}
\newcommand{\Hot}{\mathrm{Hot}}
\newcommand{\Ch}{\mathrm{Ch}}
\newcommand{\uCh}{\underline{\mathrm{Ch}}}
\newcommand{\Homext}{\mathrm{Hom\,ext}}
\newcommand{\Lotimes}{\overset{\mathbf{L}}{\otimes}}

\begin{document}

\printedpage{9}
% head: A. Grothendieck : Dérivateurs, Ch. IX, Édition des Universités de Montpellier II et Paris VII
% foot: Version 12 mars 2002

\msp{page 14}

\section*{2\quad Comparison with Quillen}

A certain number of questions arise about this functor $(\M\Sigma^{-1})^o \times (\M\Sigma^{-1}) \longrightarrow \Hot$, and its relation to the classical constructions, in particular to the theory of \textsc{Quillen}. I will now write
\[
\begin{aligned}
\Ch^{\Sigma}_{\infty}(a,b) \;&=\; \Homext_{\Hot}(a,b) && \text{or simply}\\
&\phantom{\;=\;} \Homext(a,b).
\end{aligned}
\]

1$^o$). Let $f,g : a \rightrightarrows b$ be in $\M$, then (if $\M$ is a ``model category'') \textsc{Quilllen} defines a set % [sic] Quilllen -> Quillen, three l's printed
\[
\pi_1(a,b;f,g) \qquad \text{or} \qquad \pi_1(f,g)
\]
of homotopy classes from $f$ to $g$, in such a way that $\mathrm{Hom}_{\M}(a,b)$ becomes the set of objects of a groupoid, whose arrows are the elements of $\pi_1(f,g)$, for a natural notion of composition. I suspect that one has
\[
\pi_1(f,g) \;=\; \mathrm{Hom}_{\underbrace{\Pi(\uCh^{\Sigma}_{\infty}(a,b))}_{\substack{\text{fundamental}\\ \text{groupoid}\\ \text{of } \uCh^{\Sigma}_{\infty}(a,b)}}}(f,g), % [sic] fodamental -> fondamental, missing n
\]
when $f$ and $g$ are considered as ``infinite paths'' (by inserting each $f : a \longrightarrow b$ into
\[
\cdots \underbrace{a_{-2}}_{=a} \xleftarrow{\ \mathrm{id}\ } \underbrace{a_{-1}}_{=a} \xrightarrow{\ f\ } \underbrace{a_{0}}_{=b} \xleftarrow{\ \mathrm{id}\ } \underbrace{a_{1}}_{=b} \longrightarrow \cdots\ \ )
\]

\msp{page 15}

So the definition of the $\pi_1(f,g)$ (if that is indeed it) extends in a natural way to arbitrary infinite (essentially finite) paths.

2$^o$) Suppose that $\M$ admits an object both initial and final $0_{\M}$ ($\M$ ``pointed''). I say that then all the categories $\uCh^{\Sigma}_{\infty}(a,b)$ are pointed in a natural way: namely, by the infinite path from $a$ to $b$ defined (as above) starting from $f = 0_{a,b} : a \longrightarrow b$, the zero map from $a$ to $b$. One then feels like viewing the $\uCh^{\Sigma}_{\infty}(a,b)$ as defining elements in $\Hot_{\bullet}$ (``pointed Hot''). This is finer than saying that one gives oneself a neutral class
\[
e \;\longrightarrow\; \Ch(a,b) \qquad \text{in Hot,}
\]
for the functor from $\Hot_{\bullet}$ to Hot (or to $e\backslash\Hot$) is not faithful. The problem is to define the arrows
\[
\Ch(a,b)\times\Ch(b,c) \;\longrightarrow\; \Ch(a,c)
\]
in $\Hot_{\bullet}$, and not just in Hot, and with associativity (and the units) in $\Hot_{\bullet}$ as well. If one applies the definition with $\Ch^{-}(a,b)$ and $\Ch^{+}(b,c)$, the composite of the zero elements $a \xrightarrow{\ 0\ } b$ of $\Ch^{-}(a,b)$ and $b \xleftarrow{\ \mathrm{id}\ } b \xrightarrow{\ 0\ } c$ of $\Ch^{+}(b,c)$
\printedpage{10}
% head: A. Grothendieck : Dérivateurs, Ch. IX, Édition des Universités de Montpellier II et Paris VII
% foot: Version 12 mars 2002
\msp{page 16}
is the following element of $\Ch_{\infty}(a,c)$
\[
\begin{tikzcd}[row sep=huge]
 & a_{-2} & a_{-1} \arrow[l] \arrow[r] & a_{0} \arrow[d, equal] & a_{1} \arrow[l] \arrow[r] & a_{2} & \cdots \\
{}_{\mathrm{id}}\ \cdots & a & a \arrow[l, "\mathrm{id}"'] \arrow[r, "0"'] & b & b \arrow[l, "\mathrm{id}"'] \arrow[r, "0"'] & c & \cdots_{\mathrm{id}} \arrow[l, "\mathrm{id}"']
\end{tikzcd}
\]
this is \textit{not} the zero element of $\Ch_{\infty}(a,c)$ ! But one has canonical arrows
\[
e^{-}_{a,b}\circ e^{+}_{b,c} \;\;\rightleftarrows\;\; e_{a,c} \qquad \text{in } \Ch_{\infty}(a,c)
\]
which are the identity on all the vertices except the one of index 0. There is a sorites to develop, to the effect that a pair
\[
(f : K_{\bullet} \longrightarrow L_{\bullet}, c)
\]
of a functor between pointed categories, and of a path $c : f(e_{K_{\bullet}}) \longrightarrow e_{L_{\bullet}}$, defines in Hot a canonical arrow
\[
\mathrm{hot}_{\bullet}(K_{\bullet}) \;\longrightarrow\; \mathrm{hot}_{\bullet}(L_{\bullet})
\]
with transitivity if one has
\[
(f' : L_{\bullet} \longrightarrow M_{\bullet}, c' : f'(e_{L_{\bullet}}) \longrightarrow e_{M_{\bullet}}),
\]

\msp{page 17}

whence
\[
f'' = f'\circ f : K_{\bullet} \longrightarrow M_{\bullet}, \qquad \text{and}\ \ c'' : f''(e_{K_{\bullet}}) \longrightarrow e_{M_{\bullet}}), % [sic] closing parenthesis without opening parenthesis -> (f'' = f'\circ f : ..., c'' : ...), unbalanced parenthesis
\]
defines thus
\[
\begin{tikzcd}
f'(f(e_{K_{\bullet}})) \arrow[r, "f'(c)"] \arrow[rr, bend right=12, "{c''\ \text{by definition}}"'] & f'(e_{L_{\bullet}}) \arrow[r, "c'"] & e_{M_{\bullet}}
\end{tikzcd}
\]

\noindent\llap{?\quad}So this does indeed give us pairings in $\Hot_{\bullet}$ (and also in $\mathrm{Cat\,hot}_{\bullet}$?). It is hard to believe that they could fail to be associative! Besides, $\Hot_{\bullet}$, with the ordinary cartesian product (\footnote[4]{No, one must take the product $\xi_{\bullet} \divideontimes \eta_{\bullet}$ ($\xi_{\bullet}\times\eta_{\bullet}$, with $\xi_{\bullet}\vee\eta_{\bullet}$ contracted to $e$).}), satisfies the earlier axiomatics for $\mathcal{H}$, and $\Hot_{\bullet} \longrightarrow \Hot$ is compatible with this product (\footnote[5]{No, but one gets by all the same \ldots}), and with $\pi_0$. % Reading resolved against PDF leaf 22: \divideontimes, meaning smash; see corrections.csv.

If all this works, one would thus have a theory of the $\Homext_{\bullet}$ with values in $\Hot_{\bullet}$; and probably also for the $\mathrm{CH}_{\infty}$ as objects of $\mathrm{Ind\,Hot}_{\bullet}$, and not only of $\mathrm{Ind\,Hot}$. But
\printedpage{11}
% head: A. Grothendieck : Dérivateurs, Ch. IX, Édition des Universités de Montpellier II et Paris VII
% foot: Version 12 mars 2002
in any case one cannot in general drop Cat hot, which remains important -- so one would have to have objects $\mathrm{CH}_{\infty}(a,b)$ in $\mathrm{Ind\,Cat\,hot}_{\bullet}$, to have a well-tailored setup.

The $\pi_0(C(a,b))$ will now be pointed sets, and this will be by the image of the zero map $a \xrightarrow{\ 0\ } b$, i.e. the composite
\[
a \longrightarrow 0 \longrightarrow b
\]

\msp{page 18}

which one must compose with the element deduced from the truncated path % [sic] l'elément -> l'élément, missing accent
\[
a \longrightarrow \underbrace{0}_{=a_0} \longleftarrow b \qquad \ldots
\]

One would have to show that 0 is also a neutral object of $\M\Sigma^{-1}$, i.e. that
\[
\pi_0(a,0) = \pi_0(0,a) = \{0\} \qquad \text{for every } a \in \mathrm{Ob}\,\M
\]
and interpret the zero element in $\pi_0(a,b)$ as the zero map in $\M\Sigma^{-1}$. In general, one needs conditions which ensure (without $\M$ necessarily being pointed) that
\[
\M \xrightarrow{\ \gamma\ } \M\Sigma^{-1}
\]
commutes with finite sums, with finite products (so $\gamma(\varphi_{\M}) = \varphi_{\M\Sigma^{-1}}$, resp. $\gamma(e_{\M}) = e_{\M\Sigma^{-1}}$) (\footnote[6]{to do.}).

Suppose finally that $\M$ is \textit{additive,} and even $k$-linear ($k$ a fixed ring), i.e. one has in addition a ring homomorphism
\[
k \;\longrightarrow\; \mathrm{End}(\mathrm{id}_{\M}).
\]

One must assume some compatibility of $\Sigma$ with the additive structure, e.g. that the direct sum of two arrows in $\Sigma$ is in $\Sigma$, and probably something stronger, of the ``five lemma'' kind for $\Sigma$. (But one does not necessarily assume $\M$ to be abelian, so one would have trouble formulating a five lemma?) % [sic] dû mal -> du mal, accent on du

\msp{page 19}

But perhaps here a subtler axiom is needed, e.g. (if $\underbrace{\Sigma}_{\text{suspension}}$ and $\Omega$ exist (\footnote[7]{cf. below what this means.})) that the set of adjunction morphisms $\Sigma\Omega \longrightarrow \mathrm{id}$, $\mathrm{id} \longrightarrow \Omega\Sigma$ be \texttt{[made up of]} isomorphisms, which implies that $\Omega$ resp. $\Sigma$ is fully faithful. In any case, the natural expectation here is that $\Ch^{\Sigma}_{\infty}(a,b)$ be not only in $\Hot_{\bullet}$, but even in the category $D^{\leq 0}(k\text{-Mod})$, a full subcategory of the derived category $D^{-}(k\text{-Mod})$ consisting of the $K$ such that $H^{i}(K) = 0$ if $i > 0$. One can also denote it $D_{\geq 0}(k\text{-Mod})$, interpreting it in
\printedpage{12}
% head: A. Grothendieck : Dérivateurs, Ch. IX, Édition des Universités de Montpellier II et Paris VII
% foot: Version 12 mars 2002
terms of complexes of ``chains'' rather than of cochains. Here the pairings must correspond to the operation $* = \Lotimes_k$ in $D_{\geq 0}(k\text{-Mod})$, admitting $k$ as unit object, and
\[
\pi_0(K_{\bullet}) \;\simeq\; H_0(K_{\bullet}).
\]
So one must have
\[
\Homext(a,b) \Lotimes_k \Homext(b,c) \;\longrightarrow\; \Homext(a,c),
\]
where this time Hom ext denotes the object of $D_{\geq 0}(k\text{-Mod})$, instead of just an object of Hot, or of $\Hot_{\bullet}$. One uses the functors
\[
D_{\geq 0}(k\text{-Mod}) \longrightarrow \Hot_{\bullet} \longrightarrow \Hot
\]
\msp{page 20}
to make the link with the pointed theory in $\Hot_{\bullet}$, or \texttt{[the]} general theory in Hot. If one denotes by
\[
\gamma :\ D_{\geq 0}(k\text{-Mod}) \;\longrightarrow\; \Hot_{\bullet}
\]
the first of these functors, it does not commute with the products $*$ ($\Lotimes_k$ and $\times$ respectively), but one has, for $P_{\bullet}$, $Q_{\bullet}$ varying in $D_{\geq 0}(k\text{-Mod})$, a canonical functorial arrow % [sic] variable -> variables, agreement with P, Q
\[
\gamma(P_{\bullet} \Lotimes_k Q_{\bullet}) \xleftarrow{\ \mathit{can.}\ } \gamma(P_{\bullet})\times\gamma(Q_{\bullet})
\]
``$k$-bilinear'', which shows that giving pairings
\[
P_{\bullet} * Q_{\bullet} \xrightarrow{\ c_k\ } R_{\bullet},
\]
whence
\[
\gamma(P_{\bullet} * Q_{\bullet}) \xrightarrow{\ \gamma(c_k)\ } \gamma(R_{\bullet}), \tag*{$(*)$}
\]
is in either case finer than giving oneself only
\[
\gamma(P_{\bullet}) * \gamma(Q_{\bullet}) \xrightarrow{\ c_{\bullet}\ } \gamma(R_{\bullet}), \tag*{$(*')$}
\]
seeing that a datum $(*)$ will imply a datum $(*')$ by % [sic] donne -> donnée, missing letters
\[
\begin{tikzcd}
 & \gamma(P_{\bullet} * Q_{\bullet}) \arrow[dr, "\gamma(c)"] & \\
\gamma(P_{\bullet}) * \gamma(Q_{\bullet}) \arrow[ur, "\text{can.}"] \arrow[rr, "c_{\bullet}"] & & \gamma(R_{\bullet}).
\end{tikzcd}
\]
In view of this, one does indeed expect to have a \textit{strong} theory of the pairings
\printedpage{13}
% head: A. Grothendieck : Dérivateurs, Ch. IX, Édition des Universités de Montpellier II et Paris VII
% foot: Version 12 mars 2002
\msp{page 21}
\[
\Homext_k(a,b) \Lotimes_k \Homext_k(b,c) \;\longrightarrow\; \Homext_k(a,c),
\]
from which one will deduce the weaker pairings in $\Hot_{\bullet}$ (instead of $D_{\geq 0}(k\text{-Mod})$)
\[
\Homext_{\Hot_{\bullet}}(a,b) \times \Homext_{\Hot_{\bullet}}(b,c) \;\longrightarrow\; \Homext_{\Hot_{\bullet}}(a,c).
\]

3$^o$) Returning to the case $\M$ pointed and nothing more, \textsc{Quillen} defines, for $a$, $b$ in $\M$,
\[
\pi_1(a,b) \;=\; \pi_1(a,b\,;0_a,0_b)
\]
as a bifunctor on $(\M\Sigma^{-1})^o\times(\M\Sigma^{-1})$ (with the notation of 1$^o$)). If what I put forward without proof in 1$^o$) is true, it must follow that one has a canonical isomorphism
\[
\pi_1(a,b) \;\simeq\; \pi_1(\Ch^{\Sigma}_{\infty}(a,b)),
\]
where $\Ch^{\Sigma}_{\infty}(a,b)$ is viewed as an object of $\Hot_{\bullet}$, so its $\pi_1$ is defined. One has in any case, without any reference to cofibrations and fibrations, a bifunctor with values in Groups, composed in
\[
(\M\Sigma^{-1})^o\times(\M\Sigma^{-1}) \xrightarrow{\ \ (a,b)\longmapsto \Ch^{\Sigma}_{\infty}(a,b)\ \ } \Hot_{\bullet} \xrightarrow{\ \pi_1\ } (\text{Groups}).
\]
If for fixed $a$, variable $b$,
\[
b \;\longmapsto\; \pi_1(a,b)
\]
is representable, one denotes the object representing it by $\Sigma(a)$ (\textit{suspension} -- do not confuse this $\Sigma$ with $\Sigma \subseteq \mathrm{Fl}(\M)$, which it is better now to denote by $W$, as in ``weak equivalence''). So

\msp{page 22}

\[
\mathrm{Hom}(\Sigma a,b) \;\simeq\; \pi_1(a,b),
\]
and $\Sigma a$ depends functorially on $a$, subject to existence. If it always exists, one has
\[
\Sigma :\ \M W^{-1} \;\longrightarrow\; \M W^{-1}.
\]
One defines likewise, subject to existence,
\[
\Omega :\ \M W^{-1} \;\longrightarrow\; \M W^{-1}
\]
by
\[
\mathrm{Hom}(a,\Omega b) \;\simeq\; \pi_1(a,b),
\]
and if both exist, one has the adjunction formula
\[
\mathrm{Hom}(\Sigma a,b) \;\simeq\; \mathrm{Hom}(a,\Omega b) \qquad (\simeq \pi_1\Homext_{\Hot_{\bullet}}(a,b)).
\]

\end{document}
